\documentclass[11pt]{article}
\usepackage[a4paper,margin=25mm]{geometry}
\usepackage{amsmath,amssymb,booktabs,hyperref}
\usepackage[T1]{fontenc}
\usepackage{lmodern}
\setlength{\emergencystretch}{3em}
\title{SMC$^2$ for Simulation-Assisted Toppling-Safe Pushing}
\author{Project implementation note}
\date{}
\begin{document}
\maketitle

\section{Purpose and scope}
Our agent pushes an upright cuboid toward a goal while avoiding toppling.
Geometry $G=(w,L,H)$ is observed, while physical parameters
$\theta=(m,\mu,h_c)$ (mass, floor friction, and vertical centre of mass)
are unknown. The existing CMA-ES estimator searches for low trajectory
error and returns a few good candidates. SMC$^2$ instead maintains a
weighted approximation to a Bayesian parameter posterior: alternative
physics hypotheses and their probability mass. These hypotheses configure
the internal direct-force PyBullet simulator used for action replay;
integration with imagination and the pessimistic action shield is planned. They are never appended to the SAC
policy observation.

The implementation in \texttt{src/agents/system\_identification/smc2.py} is SMC$^2$ specialized to
\emph{conditional action replay}. Each push is one observation block;
its measured start state is conditioned upon, with explicit uncertainty.
It is not a continuous latent-state tracker across pushes, and does not
model unobserved three-dimensional contact states. This matches the
existing \texttt{ActionRecord} and \texttt{InternalSim.replay}
interfaces.
CMA-ES remains available as a separate baseline.

\section{General \texorpdfstring{SMC$^2$}{SMC2} mathematics}
For a state-space model with static parameter $\theta$, latent state $x_t$,
observation $y_t$, and known action $a_t$, write
\begin{align}
 \theta &\sim p(\theta), & x_0&\sim p_\theta(x_0),\\
 x_t&\sim f_\theta(x_t\mid x_{t-1},a_t), &
 y_t&\sim g_\theta(y_t\mid x_t).
\end{align}
The target is
\begin{equation}
 p(\theta,x_{0:t}\mid y_{1:t},a_{1:t})\propto
 p(\theta)p_\theta(x_0)\prod_{r=1}^t
 f_\theta(x_r\mid x_{r-1},a_r)g_\theta(y_r\mid x_r).
\end{equation}
An outer population $\{\theta^{(j)},W_t^{(j)}\}_{j=1}^{N_\theta}$
approximates the parameter posterior. Each parameter particle owns an
inner particle filter with $N_x$ state particles and an estimated
marginal likelihood. For a bootstrap filter with resampling before
propagation, the inner step is
\begin{align}
 A_t^{(j,i)}&\sim\operatorname{Categorical}(w_{t-1}^{(j,1:N_x)}),\\
 x_t^{(j,i)}&\sim f_{\theta^{(j)}}(\cdot\mid
 x_{t-1}^{(j,A_t^{(j,i)})},a_t),\\
 v_t^{(j,i)}&=g_{\theta^{(j)}}(y_t\mid x_t^{(j,i)}),\qquad
 w_t^{(j,i)}=\frac{v_t^{(j,i)}}{\sum_k v_t^{(j,k)}},\\
 \widehat\ell_t^{(j)}&=\frac{1}{N_x}\sum_i v_t^{(j,i)},\qquad
 \widehat L_t^{(j)}=\prod_{r=1}^t\widehat\ell_r^{(j)}.
\end{align}
$A_t^{(j,i)}$ is an ancestor index, not the commanded push $a_t$.
The likelihood estimate is nonnegative and unbiased under the particle
filter construction; its logarithm is generally biased.
The outer update is
\begin{equation}
 W_t^{(j)}=\frac{W_{t-1}^{(j)}\widehat\ell_t^{(j)}}
 {\sum_k W_{t-1}^{(k)}\widehat\ell_t^{(k)}},\qquad
 \operatorname{ESS}_t=\frac{1}{\sum_j(W_t^{(j)})^2}.
\end{equation}
All computations use log likelihoods and log-sum-exp to avoid underflow.

When ESS falls below $\rho N_\theta$, resample outer particles along with
their inner filter states and likelihood estimates, reset weights to
$1/N_\theta$, and rejuvenate by particle marginal Metropolis--Hastings
(PMMH). Propose $\theta'\sim q(\cdot\mid\theta)$, run a \emph{fresh}
inner filter over all observations to obtain $\widehat L_t'$, and accept
with probability
\begin{equation}
 \alpha=\min\left(1,
 \frac{p(\theta')\widehat L_t' q(\theta\mid\theta')}
 {p(\theta)\widehat L_t q(\theta'\mid\theta)}\right).
\end{equation}
Upon rejection, retain the old parameters, inner particles, and old
random likelihood estimate. Re-estimating the denominator on rejection
would break the pseudo-marginal construction. The extended-space target
has the desired parameter posterior as its marginal; finite particle
populations and finite MCMC budgets still incur Monte Carlo error.

\section{The project's conditional replay model}
\subsection{Parameters and prior}
We work in transformed coordinates
\begin{equation}
 z=(\log m,\mu,r),\qquad r=h_c/H,\qquad
 z\sim\mathcal U([z_{\min},z_{\max}]),
\end{equation}
with bounds
\begin{equation}
 z_{\min}=(\log 0.2,0.10,0.20),\qquad
 z_{\max}=(\log 3.0,1.10,0.80).
\end{equation}
Thus the induced mass prior is log-uniform ($p(m)\propto1/m$); friction
and relative COM have uniform priors. These wider-than-training bounds
match the existing SysID search box and permit out-of-distribution
objects. This is a different prior from CMA-ES's initialization around
its training mean. No Jacobian appears in PMMH because both the target
and proposal are expressed in $z$.

\subsection{States, actions, and observations}
Record $t$ contains the measured pre-warm-up state
$b_t=(x,\psi,\dot x,\dot\psi)$, the physical commanded action
$a_t=(h,F,\tau)$, the warm-up step count, and a dense trajectory.
For this model the uncertain latent state $s_t$ is the true replay start:
\begin{equation}
 s_t\mid b_t\sim\mathcal N(b_t,S),\qquad
 S=\operatorname{diag}(\sigma_{s,1}^2,\ldots,\sigma_{s,4}^2).
\end{equation}
This conditional distribution is a modeling assumption, not a posterior
derived from a complete sensor model. The transition is independent of
the previous block: $f(s_t\mid s_{t-1},b_t)=\mathcal N(b_t,S)$.
Consequently ancestor selection in the general bootstrap filter has no
effect: each block draws fresh independent starts. This specialization
is why the code needs no inner ancestor resampling between pushes.

For each start particle, the deterministic simulator returns
$\widehat y_{t,k}=F_{G,z}(s_t,a_t)_k$, including recorded warm-up steps.
The simulator reconstructs height from planar pitch and floor contact,
and zeros other unrecorded velocity components. The observation vector is
$y_{t,k}=(x,\psi,\dot x,\dot\psi)$, in metres, radians, metres/second,
and radians/second. It is the dense physical trajectory, not the
seven-dimensional policy observation $(d,\psi,\dot x,\dot\psi,G)$.

We use $\mathcal I_t=\{k\in\{0,d,2d,\ldots\}:w_t\leq k<n_t\}$,
where $w_t$ is the recorded warm-up count and $n_t$ is the trajectory length
(default stride $d=12$ at 240 Hz). Warm-up is replayed but excluded from
scoring; its removal does not shift the sampling phase. For example,
$w_t=8$ and $d=12$ selects $12,24,\ldots$. Empty scoring windows are rejected.
The normalized likelihood is
\begin{align}
 g_z(y_t\mid s_t,a_t)&=\prod_{k\in\mathcal I_t}
 \mathcal N(y_{t,k};F_{G,z}(s_t,a_t)_k,R),\\
 \log g_z&=-\frac12\sum_{k\in\mathcal I_t}\sum_{c=1}^4
 \left[\frac{(y_{t,k,c}-\widehat y_{t,k,c})^2}{\sigma_{y,c}^2}
 +\log(2\pi\sigma_{y,c}^2)\right].
\end{align}
This answers how $p(y\mid x)$ is computed: replay the state hypothesis,
compare predicted and measured quantities, then evaluate their specified
noise density. The simulator alone does not supply an observation density.
Pitch residuals are ordinary differences within the upright regime;
this model is not intended for wrapped rotations after a full topple.

Default start standard deviations are
$(0.01,0.002,0.01,0.02)$ in the channel units above. Observation standard
deviations default to the inverse square roots of \texttt{LOSS\_SCALES}:
\begin{equation}
 \sigma_{y,c}=1/\sqrt{\texttt{LOSS\_SCALES}[c]},\qquad
 \sigma_y\approx(0.00356,0.000470,0.0204,0.0437).
\end{equation}
Both are configurable. Replay error is not automatically independent
Gaussian sensor noise: these defaults are provisional. Subsampling
reduces, but does not prove the absence of, temporal correlation.
The implementation rejects empty, nonfinite, mismatched, or incorrectly
timed trajectories rather than silently truncating them.

\subsection{Likelihood estimator and parameter update}
For each outer particle, draw $s_t^{(j,i)}\sim\mathcal N(b_t,S)$ and set
\begin{equation}
 \widehat\ell_t^{(j)}=\frac1{N_x}\sum_i
 g_{z^{(j)}}(y_t\mid s_t^{(j,i)},a_t),\qquad
 w_t^{(j,i)}\propto g_{z^{(j)}}(y_t\mid s_t^{(j,i)},a_t).
\end{equation}
This is unbiased for the conditional block likelihood integral. With
independent random draws across blocks, the product is unbiased for
\begin{equation}
 L_T(z)=\prod_{t=1}^T\int g_z(y_t\mid s_t,a_t)
 \mathcal N(s_t;b_t,S)\,ds_t.
\end{equation}
We target $p(z\mid y_{1:T},b_{1:T},a_{1:T})\propto p(z)L_T(z)$
under this conditional model, not the full likelihood of the measured
starts and all dense trajectories together. Correlation between the
previous endpoint and the next measured start is not modeled.
Zero start variance is allowed for deterministic replay experiments;
then inner particles coincide and offer no additional uncertainty.

Outer resampling is systematic. With \texttt{proposal="fixed"}, PMMH uses
$z'=z+\epsilon$ with independent
$\epsilon_c\sim\mathcal N(0,[\eta(z_{\max,c}-z_{\min,c})]^2)$,
default $\eta=0.08$. Proposals outside the bounds are rejected without
clipping. Inside the bounds the symmetric proposal and uniform prior
cancel, so $\log\alpha=\min(0,\log\widehat L_T'-\log\widehat L_T)$.
Each accepted proposal replaces the complete current filter state.

The default \texttt{proposal="covariance"} instead uses the weighted
particle covariance in prior-scaled coordinates, measured before resampling.
Its shape is regularized and normalized to have trace three, keeping the same
expected squared step length as the fixed proposal in those coordinates.
A collapsed population falls back to the fixed shape. The transform is frozen
throughout each move stage, so the Gaussian proposal remains symmetric.
This population-based adaptation does not guarantee exact posterior samples
from a finite population.

Optional \texttt{temper\_updates=True} introduces each new push through
likelihood powers from zero to one while keeping earlier evidence at full
strength. Conditional ESS determines the increments. This mode requires
\texttt{start\_std=(0,0,0,0)} and \texttt{n\_state=1}; it rejects uncertain
starts and \texttt{ess\_fraction=1}. Each intermediate resampling stage uses
the configured proposal and move count, so gradual updates can cost more
replays. The default remains a full likelihood update.

\section{Using the implementation}
\begin{verbatim}
from src.agents.system_identification.smc2 import SMC2SysID
from src.simulation.modelling.actions import PushAction

history = []  # persists across episodes for one object
with SMC2SysID(geometry=(w, L, H), n_theta=64, n_state=8,
              seed=7) as estimator:
    record = target.push(PushAction(h=h, F=F, tau=tau), warmup_steps=8)
    # Select upright records under a declared data-collection protocol.
    history.append(record)
    result = estimator.update(history)
    z, weights = estimator.particles, estimator.weights
\end{verbatim}
The record uses physical actions, not normalized gym actions. Repeated
\texttt{update(history)} calls process only appended records. Prefix
edits or shortened histories raise an error; construct a new estimator
for a new object or revised data. Empty history returns prior hypotheses.
The estimator snapshots records so later caller edits cannot silently
change the observations used in PMMH. Its context manager closes its
owned PyBullet client. A supplied simulator remains caller-owned.

\texttt{result.hypotheses} provides the existing \texttt{SysIDResult}
interface: duplicate parameter particles have their weights combined,
and the top $K$ unique particles have weights renormalized over that
subset. \texttt{point} is the particle with greatest aggregated mass,
not a continuous MAP or minimum-error estimate. After resampling,
weights can tie, so this point is a coarse summary.
\texttt{loss} is negative estimated log likelihood, possibly negative,
and is not comparable to the CMA-ES squared-error loss.

Use the full weighted population for uncertainty. First convert to physical
parameters $\theta_j=(\exp z_{j,1},z_{j,2},Hz_{j,3})$. Then compute
\begin{equation}
 \bar\theta=\sum_j W_j\theta_j,\qquad
 \operatorname{Cov}(\theta)=\sum_jW_j
 (\theta_j-\bar\theta)(\theta_j-\bar\theta)^\top.
\end{equation}
Top-$K$ truncation can omit important low-probability dangerous dynamics;
a shield using only that subset is not a posterior safety guarantee.
The new estimator does not implement the future SAC/Dyna/shield stages.

Run the small independent-episode example from the repository root:
\begin{verbatim}
.venv/bin/python -m \
  src.evaluation.system_identification.demos.smc2_demo \
  --n-theta 64 --n-state 8
.venv/bin/python -m pytest tests/test_smc2.py
mkdir -p artifacts/smc2-doc
pdflatex -halt-on-error -interaction=nonstopmode \
  -output-directory=artifacts/smc2-doc docs/smc2.tex
\end{verbatim}
The demo reports weighted ensemble means and standard deviations, ESS
before resampling, and PMMH acceptance counts. True parameters are used
only to configure the target and for evaluation output.

\section{Limitations, cost, and validation}
A new push costs $N_\theta N_x$ complete replays. A rejuvenation event
with $M$ PMMH steps after $T$ records costs up to
$M N_\theta N_x T$ replays. Defaults are $N_\theta=64$, $N_x=8$,
$M=2$, and $\rho=0.5$. The code stores the most recent block's inner
particles, cumulative likelihoods, and the recorded history; it does
not store every simulated trajectory. Adaptive inner population sizes and
parallel replay are not implemented. The cost above describes full updates;
the optional known-start gradual updates can add resampling/move stages.

Dense evidence can collapse a small outer population. Low inner
likelihood precision can cause sticky PMMH chains. Examine ESS,
acceptance, seed-to-seed stability, particle diversity, and posterior
predictive coverage; increase both budgets as needed. More particles
cannot fix an incorrect observation model. Calibrate noise and temporal
subsampling on held-out replay trajectories before drawing scientific
conclusions about credible intervals or safety.

Sliding excitation often identifies mass and friction better than COM
height. Near-tipping pushes can add COM information but replay is
bistable; complete topples are outside this model's intended regime.
Data collection and censoring rules must be reported: discarding outcomes
based on observed tilt can itself affect inference, and this version
does not include a selection/censoring likelihood. The arm-versus-direct-
force discrepancy will also require a revised model once the KUKA
executor is integrated.

Tests compare sequential weights with exact Gaussian Bayes updates,
check the inner evidence and state mean against an analytic Gaussian
convolution, check rejuvenation against a known posterior, verify
rejection retains the random likelihood, and exercise real PyBullet
replay and resource cleanup. These establish algorithmic correctness
for the tested cases, not empirical parameter recovery on all objects.

\begin{thebibliography}{9}
\bibitem{chopin} N. Chopin, P. E. Jacob, and O. Papaspiliopoulos.
SMC$^2$: an efficient algorithm for sequential analysis of state space
models. \emph{Journal of the Royal Statistical Society: Series B},
75(3):397--426, 2013. \url{https://doi.org/10.1111/j.1467-9868.2012.01046.x}.
\bibitem{andrieu} C. Andrieu, A. Doucet, and R. Holenstein.
Particle Markov chain Monte Carlo methods.
\emph{Journal of the Royal Statistical Society: Series B},
72(3):269--342, 2010. \url{https://doi.org/10.1111/j.1467-9868.2009.00736.x}.
\end{thebibliography}
\end{document}
